\documentclass[11pt,a4paper]{article}
\usepackage[margin=2.5cm]{geometry}
\usepackage[T1]{fontenc}
\usepackage{lmodern}
\usepackage{booktabs}
\usepackage{enumitem}
\usepackage{xcolor}
\usepackage{hyperref}
\hypersetup{colorlinks=true,linkcolor=black,urlcolor=blue!60!black}
\setlist{nosep}
\setlength{\parskip}{4pt}
\setlength{\parindent}{0pt}

\title{NetRadiant-custom: Modern Lighting Pipeline Plan\\
\large Direct PBR, shadow mapping, and SH indirect lighting for a Silent Hill 1 style game}
\author{Davis Wen}
\date{16 September 2026}

\begin{document}
\maketitle

\section{Goal}

Extend NetRadiant-custom and its build pipeline so a level designer can author, preview, and compile
maps with a glTF-style metallic--roughness material model, physically based direct lighting with
shadow maps, and baked indirect lighting stored as spherical harmonics (SH). The target game is a
Silent Hill 1 style third-person title running on an ioquake3-derived engine: a fog-limited town of
watertight streets under a low virtual sky, with watertight interiors.

The plan is staged so that every stage ends with something a level designer can look at, and no
stage requires undoing the previous one.

\section{Guiding decisions}

\begin{description}[style=nextline]
\item[Keep q3map2's BSP and VIS stages untouched.]
  CSG, the BSP tree, portals, leak detection, t-junctions, patch tessellation, fog volumes,
  collision lumps, and the file formats are exactly what we do not want to rewrite.
\item[Replace the light stage with a separate tool.]
  The existing \texttt{sh-baker} (Embree, tinygltf) gains a BSP reader and becomes the light stage.
  q3map2's \texttt{-light} is used only to allocate lightmap UVs, or dropped once the baker atlases
  on its own.
\item[No BSP format change for thick walls.]
  The BSP brush and brushside lumps keep every plane of every solid brush, including void-facing
  planes whose draw surfaces q3map2 strips. The engine reconstructs closed convex hulls from those
  lumps at load and uses them as front-face-culled shadow casters.
\item[Editor preview of direct lighting is the fast iteration loop.]
  Editor preview of baked SH is optional and last, because it needs lightmap UVs on source
  geometry, which is a new subsystem.
\item[Lock units and colour spaces in Stage~1.]
  Every later comparison against the engine depends on light units, exposure, and sRGB handling
  being final in the benchmark map.
\end{description}

\section{Current state of the editor (relevant facts)}

\begin{itemize}
\item Gamepacks are data: \texttt{gamepacks/games/*.game} plus a matching directory. The editor
  scans them at start; no code change is needed to add a game.
\item Material parsing is compiled into \texttt{plugins/shaders}. Three hard-coded languages exist
  (Quake~3, Doom~3, Quake~4), selected by the \texttt{shaders} attribute of the \texttt{.game} file.
  The Doom~3/Quake~4 path already resolves diffuse, bump, and specular textures per material.
\item The camera has a lighting draw mode that renders one additive GLSL pass per light
  (\texttt{radiant/renderstate.cpp}, \texttt{GLSLBumpProgram}) using map light entities as omni
  lights with attenuation textures. Brush and patch vertices already carry tangent and bitangent.
\item Rendering goes straight to the 8-bit window buffer; there is no framebuffer object in the
  camera path. The context is created by \texttt{QOpenGLWidget} with no explicit version request.
\item The editor has no notion of lightmap UVs; q3map2's light stage allocates them.
\item The build menu is a per-gamepack XML list of command lines; each q3map2 stage is a separate
  invocation that reads and writes the BSP.
\end{itemize}

\section{Stage 1: Direct light and PBR materials}

\textbf{Outcome.} A gamepack, a material format, light entity definitions, and a lit editor preview
good enough to author the first benchmark map with final storage formats.

\subsection{1.1 Material format and parser}
\begin{itemize}
\item Define a text material format modelled on the glTF metallic--roughness schema, brace-compatible
  with the Doom~3 parser: \texttt{basecolor}, \texttt{normal}, \texttt{metallicroughness},
  \texttt{occlusion}, \texttt{emissive}, plus \texttt{metallicFactor}, \texttt{roughnessFactor},
  \texttt{emissiveFactor}, \texttt{emissiveStrength}, alpha mode and cutoff, double-sided.
\item Keep \texttt{qer\_editorimage}, \texttt{surfaceparm}, and \texttt{q3map\_} keys so q3map2 and
  the editor's brush classification keep working. Decide whether the baker reads the same files.
\item Add a fourth shader language: new \texttt{parsePBR} on \texttt{ShaderTemplate}, a new
  \texttt{SHADERLANGUAGE\_} value, a module in \texttt{plugins/shaders/plugin.cpp} setting the
  directory and extension, and a \texttt{shaders="pbr"} value for the \texttt{.game} file.
\item Extend \texttt{IShader} in \texttt{include/ishaders.h} with the extra textures and factors.
\item Colour space: base colour and emissive are sRGB and decoded in the shader; normal, metallic--
  roughness, and occlusion are linear. Check the texture preference paths (gamma, compression) so the
  loaders do not alter these.
\end{itemize}

\subsection{1.2 Light entities and units}
\begin{itemize}
\item Entity definitions for point, spot, and sun lights with the keys the engine will read:
  colour, intensity in physical units, radius, cone angles, and shadow flags.
\item The sun lives on worldspawn or a dedicated entity, so Stage~2 can read it without a format
  change. Direction and colour also feed \texttt{q3map\_sun} for compatibility.
\item Plumb the new keys through the entity plugin into \texttt{RendererLight} (origin, colour,
  radius exist; add direction, cone, intensity).
\end{itemize}

\subsection{1.3 Float framebuffer and tonemap}
\begin{itemize}
\item Add a half-float colour target plus depth in the camera path, with a resolve pass that applies
  exposure and a tonemap. Physical light units clip in an 8-bit buffer before tonemapping can happen.
\item Pulled forward from Stage~2 because shadow maps need framebuffer objects anyway.
\end{itemize}

\subsection{1.4 Lighting program}
\begin{itemize}
\item New GLSL program replacing Blinn--Phong: GGX distribution, Smith geometry, Schlick Fresnel,
  Disney diffuse, ported from \texttt{sh-renderer/glsl/radiance.frag}.
\item Keep the editor's per-light additive pass structure. Do not port Forward+ tiling; the editor
  already culls lights per surface on the CPU.
\item Base pass draws emissive plus a constant ambient. Point and spot lights accumulate on top.
\end{itemize}

\subsection{1.5 Gamepack and benchmark map}
\begin{itemize}
\item Gamepack: \texttt{.game} file, entity definitions, default build menu, shader list.
\item Benchmark map: one street block with two building interiors, a low sky ceiling, a fog volume,
  a sun, several point and spot lights, and materials exercising every map type.
\item Exit criterion: the map looks the same in the editor lighting mode and in the engine's
  direct-light-only path, within exposure tolerance.
\end{itemize}

\section{Stage 2: Shadow mapping}

\textbf{Outcome.} Sun and spot shadows in the editor preview, cast by closed brush volumes, and a
proof in the engine that hull reconstruction from the BSP gives the same result.

\subsection{2.1 Caster pass}
\begin{itemize}
\item Depth-only program drawing all brush faces with front-face culling. Because every brush is a
  closed convex volume in the editor, the far side of each wall becomes the caster and the effective
  bias is the wall thickness (second-depth technique). No hand-tuned bias per map.
\item Cache the shadow maps and invalidate on scene change or light change; the editor scene changes
  rarely, so per-frame cost is only sampling.
\end{itemize}

\subsection{2.2 Sun cascades, then spot atlas}
\begin{itemize}
\item Cascaded shadow maps for the sun first, since a low sky ceiling makes the thin-wall question
  hardest there. Then the spot light shadow atlas.
\item Port \texttt{ComputeShadow}, \texttt{ComputeSunShadow}, and the Poisson PCF from
  \texttt{radiance.frag} unchanged once the samplers and cascade uniforms exist.
\end{itemize}

\subsection{2.3 Thin receivers}
\begin{itemize}
\item Patches and models are thin in the editor too. Use normal-offset bias for them, or make them
  receivers only.
\end{itemize}

\subsection{2.4 Engine-side hull reconstruction (validation)}
\begin{itemize}
\item In ioquake3, rebuild each collision brush's convex hull from the brush and brushside lumps at
  load, and use the hulls as the front-face-culled caster mesh. Draw surfaces remain the receivers.
\item Known gaps: brushes with non-solid shaders are dropped from the collision lumps. If such a
  brush must cast, mark its material so the compiler keeps it, or emit it separately. q3map2's bevel
  planes are redundant on the polytope and do not disturb the intersection.
\item Exit criterion: the benchmark map's shadows match between editor and engine.
\end{itemize}

\section{Stage 3: Lightmap UVs, SH indirect, and the compiled format}

\textbf{Outcome.} Baked SH indirect lighting in the engine via a light stage replacement, with no
change to the BSP format.

\subsection{3.1 Baker as the light stage}
\begin{itemize}
\item Add a BSP reader to \texttt{sh-baker}: draw surfaces, lightmap UVs, materials, and brush hulls
  for casters. Entities for lights and sun.
\item Pipeline: \texttt{q3map2 -meta}, \texttt{q3map2 -vis}, \texttt{q3map2 -light -fast} for UV
  allocation only, then \texttt{sh\_baker} in place of the real light stage.
\item Later, optionally atlas in the baker and write lightmap coordinates into the draw surfaces
  directly, dropping the \texttt{-light} step.
\end{itemize}

\subsection{3.2 Output and engine loading}
\begin{itemize}
\item External lightmap images (already supported by the engine) for the L0 band, plus a side file
  with the three packed SH textures matching \texttt{GetLightmapTexel} in \texttt{radiance.frag}.
\item Engine loads the side file and runs the indirect BRDF. No BSP lump changes.
\end{itemize}

\subsection{3.3 Editor preview of baked SH (optional, last)}
\begin{itemize}
\item Requires lightmap UVs on source brush faces and patches that match the compiled atlas, either
  by matching surfaces back to the BSP or by an editor-side atlas (xatlas or similar).
\item Only worth doing if designers need indirect preview inside the editor; the engine and the
  existing \texttt{sh-renderer} viewer cover it otherwise.
\end{itemize}

\section{World-building conventions for the town}

These are authoring rules, not code, but they determine whether the visibility system works.

\begin{itemize}
\item \textbf{Low sky ceiling just above the camera.} VIS treats sky as opaque, so the town becomes a
  network of corridors walled by buildings, and the PVS behaves like an indoor level.
\item \textbf{Building shells are structural brushes.} Only structural brushes block VIS. Ornament is
  detail. Props, pipes, and clutter are models.
\item \textbf{Fog colour equals sky colour.} Buildings end at the ceiling; with matched colours both
  fade to the same value and the seam disappears. Far plane set at fog opacity distance.
\item \textbf{Hint brushes across long streets} at intervals so far clusters drop out; area portals in
  doors between street and interiors.
\item \textbf{Clamp camera pitch} so the ceiling is never inspected from below.
\item \textbf{Set worldspawn \texttt{\_blocksize}} to zero or a large value for the flat town to
  avoid pointless axial splits.
\end{itemize}

\section{Risks}

\begin{tabular}{@{}p{0.34\textwidth}p{0.62\textwidth}@{}}
\toprule
Risk & Mitigation \\
\midrule
Light units or colour space differ between editor and engine &
  Fixed in Stage~1.5 by comparing the benchmark map; treat as a blocking exit criterion. \\
Texture loader preferences alter colour data &
  Audit gamma and compression paths before choosing where linearisation happens. \\
Non-solid brushes needed as casters are dropped from collision lumps &
  Material key to keep them, or a small editor-side caster export. \\
Long straight streets defeat the PVS &
  Fog far plane plus hint brushes; accept that the fog is the real culler there. \\
Editor SH preview is expensive &
  Scoped optional and last; not on the critical path. \\
Qt GL context lacks needed features for FBO and half-float &
  Request an explicit 3.3 compatibility profile in \texttt{QSurfaceFormat}; the legacy code keeps working. \\
\bottomrule
\end{tabular}

\section{Dependency summary}

\begin{tabular}{@{}llp{0.55\textwidth}@{}}
\toprule
Item & Depends on & Note \\
\midrule
1.4 Lighting program & 1.1, 1.2, 1.3 & Needs materials, light keys, and a float target. \\
1.5 Benchmark map & 1.1--1.4 & Locks units and formats. \\
2.1 Caster pass & 1.3 & Reuses FBO plumbing. \\
2.2 Cascades & 1.2 & Needs the sun definition. \\
2.4 Engine hulls & 1.5 & Needs the benchmark map to compare against. \\
3.1 Baker & 1.1, 1.5 & Reads the material format; validated on the benchmark map. \\
3.2 Engine SH loading & 3.1, 2.4 & Direct and indirect combine in the engine's radiance shader. \\
3.3 Editor SH preview & 3.2 & Optional. \\
\bottomrule
\end{tabular}

\end{document}
